\subsection{Wireframe, Mockup và Prototype}
\label{sec:wireframes-prototypes}

\begin{guide}
\guidepart{Nội dung}
\begin{compactitem}
    \item Low-Fidelity (Lo-Fi): bản phác thảo nhanh bố cục các màn hình chính, chưa có màu sắc và hình ảnh thật.
    \item High-Fidelity (Hi-Fi): giao diện hoàn chỉnh với màu sắc, font chữ và hình ảnh theo design system ở mục~\ref{sec:design-system}.
    \item Interactive Prototype: bản mô phỏng có thể bấm được (clickable), kèm link để xem.
\end{compactitem}
\guidepart{Câu hỏi cần trả lời}
\begin{compactitem}
    \item Những màn hình nào được thiết kế, và chúng bao phủ các luồng nào ở mục~\ref{sec:user-flows}?
    \item Thiết kế đã thay đổi gì từ Lo-Fi sang Hi-Fi, và vì sao (dựa trên góp ý hay kết quả kiểm thử)?
    \item Prototype được dùng để kiểm thử với ai, và kết quả ra sao?
\end{compactitem}
\guidepart{Hình/bảng nên có}
Ảnh Lo-Fi và Hi-Fi đặt cạnh nhau cho từng màn hình chính, link tới prototype.
\end{guide}
